\section{春季星空與天體}
\begin{itemize}
    \item 春季大弧線：北斗七星的玉衡、開陽、搖光（柄杓尾端三顆星）、牧夫座大角星、室女座角宿、烏鴉座
    \item 春季大三角：獅子座的五帝座一、室女座的角宿一及牧夫座的大角星
    \item 春季大鑽石：春季大三角加上獵犬座的常陳一
\end{itemize}

\subsection{大熊座 Ursa Major}
\begin{enumerate}
    \item 最亮星:大熊座\(\alpha\)星（天樞）、\(\epsilon\)星（玉衡），均為 1.8 等
    \item 北斗七星所在
    \item 面積排名第三
    \item 大熊座\(\zeta\)星（古代中國稱為開陽，北斗六）：位於大熊座中間的\(\zeta\)星與它附近的輔星構成了一對裸眼可見的雙星。大熊座\(\zeta\)星本身又是一個雙星系統，它是1650年發現的，是望遠鏡時代最早發現的雙星
\end{enumerate}
大熊座中的深空天體：

\begin{center}
\begin{longtable}{c | c | c | c | c | c}
    % \hline
    編號 & 種類 & 別名 & 視尺寸 & RA & Dec \\
    \hline
    \endfirsthead

    % \hline
    % 編號 & 種類 & 別名 & 視尺寸 & 拍攝技巧 & 曝光時間 \\
    % \hline
    % \endhead

    % \hline
    % \multicolumn{6}{r}{接下頁} \\
    \endfoot

    % \hline
    \endlastfoot

    M81/NGC 3031 & 螺旋星系 & 波德星系 & \(\ang{;21;38}\times\ang{;11;15}\) & \(9^\text{h }55^\text{m }32^\text{s}\) & \(+\ang{69;3;55}\) \\
    M82/NGC 3034 & 不規則星系 & 雪茄星系 & \(\ang{;11;}\times\ang{;5;6}\) & \(9^\text{h }55^\text{m }53^\text{s}\) & \(+\ang{69;40;45}\) \\
    M97/NGC 3587 & 行星狀星雲 & 貓頭鷹星雲 & \(\ang{;2;50}\) & \(11^\text{h }14^\text{m }48^\text{s}\) & \(+\ang{55;1;8}\) \\
    M101/NGC 5457 & 螺旋星系 & 風車星系 & \(\ang{;24;}\) & \(14^\text{h }03^\text{m }12^\text{s}\) & \(+\ang{54;20;56}\) \\
    M108/NGC 3556 & 螺旋星系 & - & \(\ang{;4;}\times\ang{;1;20}\) & \(11^\text{h }11^\text{m }30^\text{s}\) & \(+\ang{55;40;26}\) \\
    M109/NGC 3992 & 螺旋星系 & - & \(\ang{;8;}\times\ang{;5;20}\) & \(11^\text{h }57^\text{m }36^\text{s}\) & \(+\ang{53;22;28}\) \\
    NGC 2481 & 棒旋星系 & - & \(\ang{;6;54}\times\ang{;3;18}\) & \(9^\text{h }22^\text{m }02^\text{s}\) & \(+\ang{50;58;35}\) \\
\end{longtable}
\end{center}

\subsection{牧夫座 Bootes}
\begin{enumerate}
    \item 最亮星：大角星（牧夫座\(\alpha\)），橙巨星，視星等-0.05為全天第四亮星
    \item 意思是牧羊人或農夫
    \item 象限儀座流星雨的輻射點位於牧夫座中
\end{enumerate}
牧夫座中的深空天體：

\begin{center}
\begin{longtable}{c | c | c | c | c | c}
    % \hline
    編號 & 種類 & 別名 & 視尺寸 & RA & Dec \\
    \hline
    \endfirsthead

    % \hline
    % 編號 & 種類 & 別名 & 視尺寸 & 拍攝技巧 & 曝光時間 \\
    % \hline
    % \endhead

    % \hline
    % \multicolumn{6}{r}{接下頁} \\
    \endfoot

    % \hline
    \endlastfoot

    NGC 5466 & 球狀星團 & - & \(\ang{;6;36}\) & \(14^\text{h }05^\text{m }26^\text{s}\) & \(+\ang{28;32;4}\) \\
\end{longtable}
\end{center}

\subsection{室女座 Virgo}
\begin{enumerate}
    \item 最大的黃道星座，全部星座排名第二大
    \item 最亮星：角宿一（室女座\(\alpha\)星），視星等為0.98
    \item 室女座流星雨
\end{enumerate}
室女座中的深空天體：

\begin{center}
\begin{longtable}{c | c | c | c | c | c}
    % \hline
    編號 & 種類 & 別名 & 視尺寸 & RA & Dec \\
    \hline
    \endfirsthead

    % \hline
    % 編號 & 種類 & 別名 & 視尺寸 & 拍攝技巧 & 曝光時間 \\
    % \hline
    % \endhead

    % \hline
    % \multicolumn{6}{r}{接下頁} \\
    \endfoot

    % \hline
    \endlastfoot

    M58/NGC 4579 & 螺旋星系 & - & \(\ang{;5;}\times\ang{;3;50}\) & \(12^\text{h }37^\text{m }44^\text{s}\) & \(+\ang{11;49;5}\) \\
    M59/NGC 4621 & 橢圓星系 & - & \(\ang{;5;}\times\ang{;3;12}\) & \(12^\text{h }42^\text{m }02^\text{s}\) & \(+\ang{11;38;49}\) \\
    M61/NGC 4303 & 螺旋星系 & Swelling Spiral & \(\ang{;6;54}\) & \(12^\text{h }21^\text{m }55^\text{s}\) & \(+\ang{4;28;25}\) \\
    M89/NGC 4552 & 橢圓星系 & - & \(\ang{;8;}\) & \(12^\text{h }35^\text{m }39^\text{s}\) & \(+\ang{12;33;22}\) \\
    M90/NGC 4569 & 螺旋星系 & - & \(\ang{;9;6}\times\ang{;3;50}\) & \(12^\text{h }36^\text{m }50^\text{s}\) & \(+\ang{13;9;46}\) \\
    M104/NGC 4594 & 螺旋星系 & 氈帽星系 & \(\ang{;8;27}\times\ang{;4;55}\) & \(12^\text{h }39^\text{m }59^\text{s}\) & \(-\ang{11;37;23}\) \\
    M49/NGC 4472 & 橢圓星系 & - & \(\ang{;10;}\) & \(12^\text{h }29^\text{m }47^\text{s}\) & \(+\ang{8;0;1}\) \\
    M60/NGC 4649 & 橢圓星系 & - & \(\ang{;6;47}\times\ang{;5;27}\) & \(12^\text{h }43^\text{m }40^\text{s}\) & \(+\ang{11;33;9}\) \\
    M84/NGC 4374 & 橢圓星系 & （屬於馬卡萊恩長鏈） & \(\ang{;5;}\times\ang{;4;24}\) & \(12^\text{h }25^\text{m }04^\text{s}\) & \(+\ang{12;53;13}\) \\
    M86/NGC 4406 & 橢圓星系 & （屬於馬卡萊恩長鏈） & \(\ang{;11;30}\times\ang{;8;25}\) & \(12^\text{h }26^\text{m }12^\text{s}\) & \(+\ang{12;56;46}\) \\
    M87 & 橢圓星系 & 室女座A星系 & \(\ang{;7;}\) & \(12^\text{h }30^\text{m }48^\text{s}\) & \(+\ang{12;23;27}\)
\end{longtable}
\end{center}

\subsection{獅子座 Leo}
\begin{enumerate}
    \item 黃道星座之一
    \item 最亮星：軒轅十四（獅子座\(\alpha\)），視星等為1.35，藍白色恆星，含義是小皇帝
    \item 獅子前半身：鐮刀/問號 由南向北
    \item 五帝座一（獅子座\(\beta\)），視星等2.14等的白色恆星，是獅子座的第二亮星，英文含義是獅子的尾巴
    \item 軒轅十二（獅子座 \(\gamma\)），視星等1.98的聯星，呈現橘紅色和黃或黃綠色。字義是額頭，但實際上代表的是獅子的鬃毛
    \item 獅子座的設立已經數千年的曆史。現在普遍認同的說法是在4000多年前的古埃及，每年仲夏節太陽移到獅子座天區時，尼羅河的河谷就有大量獅子從沙漠中聚集乘涼喝水，獅子座因此得名。
\end{enumerate}
獅子座中的深空天體：

\begin{center}
\begin{longtable}{c | c | c | c | c | c}
    % \hline
    編號 & 種類 & 別名 & 視尺寸 & RA & Dec \\
    \hline
    \endfirsthead

    % \hline
    % 編號 & 種類 & 別名 & 視尺寸 & 拍攝技巧 & 曝光時間 \\
    % \hline
    % \endhead

    % \hline
    % \multicolumn{6}{r}{接下頁} \\
    \endfoot

    % \hline
    \endlastfoot

    \makecell{M65/NGC 3623\\M66/NGC 3627\\NGC 3628} & \makecell{螺旋星系\\螺旋星系\\棒旋星系} & \makecell{獅子座三胞胎\\Leo Triplet} & \(\ang{;50;}\times\ang{;40;}\) & \makecell{\(11^\text{h }18^\text{m }56^\text{s}\)\\\(11^\text{h }20^\text{m }15^\text{s}\)\\\(11^\text{h }20^\text{m }17^\text{s}\)} & \makecell{\(+\ang{13;5;32}\)\\\(+\ang{12;59;30}\)\\\(+\ang{13;35;23}\)} \\  
    M95/NGC 3351 & 棒旋星系 & - & \(\ang{;7;12}\) & \(10^\text{h }43^\text{m }58^\text{s}\) & \(+\ang{11;42;13}\) \\
    M96/NGC 3368 & 中間螺旋星系 & - & \(\ang{;8;18}\) & \(10^\text{h }46^\text{m }45^\text{s}\) & \(+\ang{11;49;11}\) \\
    M105/NGC 3379 & 橢圓星系 & - & \(\ang{;5;}\times\ang{;4;15}\) & \(10^\text{h }47^\text{m }49^\text{s}\) & \(+\ang{12;34;53}\) \\
\end{longtable}
\end{center}

\subsection{獵犬座 Canes Venatici}
\begin{enumerate}
    \item 最亮星:常陳一（獵犬座α星），視星等2.90等
    \item 次亮星：勢四，視星等3.83的橙巨星
    \item 代表牧夫座牽的兩條狗是三個代表狗的星座之一
\end{enumerate}
獵犬座中的深空天體：

\begin{center}
\begin{longtable}{c | c | c | c | c | c}
    % \hline
    編號 & 種類 & 別名 & 視尺寸 & RA & Dec \\
    \hline
    \endfirsthead

    % \hline
    % 編號 & 種類 & 別名 & 視尺寸 & 拍攝技巧 & 曝光時間 \\
    % \hline
    % \endhead

    % \hline
    % \multicolumn{6}{r}{接下頁} \\
    \endfoot

    % \hline
    \endlastfoot

    M3/NGC 5272 & 球狀星團 & - & \(\ang{;16;12}\) & \(13^\text{h }42^\text{m }11^\text{s}\) & \(+\ang{28;22;31}\) \\
    M51/NGC 5194 & 中間螺旋星系 & 渦狀星系 & \(\ang{;13;43}\times\ang{;11;40}\) & \(13^\text{h }29^\text{m }53^\text{s}\) & \(+\ang{47;11;42}\) \\
    M94/NGC 4736 & 螺旋星系 & 鱷魚眼星系 & \(\ang{;11;12}\times\ang{;9;6}\) & \(12^\text{h }50^\text{m }53^\text{s}\) & \(+\ang{41;7;13}\) \\
    M106/NGC 4258 & 螺旋星系 & - & \(\ang{;17;}\times\ang{;7;15}\) & \(12^\text{h }18^\text{m }57^\text{s}\) & \(+\ang{47;18;14}\)
\end{longtable}
\end{center}

\subsection{小獅座 Leo Minor}
\begin{enumerate}
    \item 是北天暗淡小星座，名稱對應獅子座
\end{enumerate}
小獅座中的深空天體：

\begin{center}
\begin{longtable}{c | c | c | c | c | c}
    % \hline
    編號 & 種類 & 別名 & 視尺寸 & RA & Dec \\
    \hline
    \endfirsthead

    % \hline
    % 編號 & 種類 & 別名 & 視尺寸 & 拍攝技巧 & 曝光時間 \\
    % \hline
    % \endhead

    % \hline
    % \multicolumn{6}{r}{接下頁} \\
    \endfoot

    % \hline
    \endlastfoot

    NGC 3432 & 棒旋星系 & 織針星系 & \(\ang{;7;30}\times\ang{;2;}\) & \(10^\text{h }52^\text{m }31^\text{s}\) & \(+\ang{36;37;7}\) \\
\end{longtable}
\end{center}

\subsection{后髮座 Coma Berenices}
\begin{enumerate}
    \item 最亮星：周鼎一（后髮座\(\beta\)），視星等為4.26
    \item 以前代表獅子尾毛處於銀河系北極方向，河外星系不會受到銀河系平面星際物質的遮掩可以觀測到大量壯觀的深空天體，包括室女座星系團和后髮座星系團等
    \item 傳說：是古代昔蘭尼的公主，埃及法老托勒密三世的王后。約在公元前243年的時候，托勒密三世遠征敘利亞。她向神祈禱保佑丈夫平安歸來，把頭髮剪下來奉獻在女神阿佛羅狄忒的神廟。傳說第二天，她的頭髮不見了。宮廷里的人說是被女神攝到了天上，成為后髮座。
\end{enumerate}
后髮座中的深空天體：

\begin{center}
\begin{longtable}{c | c | c | c | c | c}
    % \hline
    編號 & 種類 & 別名 & 視尺寸 & RA & Dec \\
    \hline
    \endfirsthead

    % \hline
    % 編號 & 種類 & 別名 & 視尺寸 & 拍攝技巧 & 曝光時間 \\
    % \hline
    % \endhead

    % \hline
    % \multicolumn{6}{r}{接下頁} \\
    \endfoot

    % \hline
    \endlastfoot

    M53/NGC 5024 & 球狀星團 & - & \(\ang{;9;}\) & \(13^\text{h }12^\text{m }55^\text{s}\) & \(+\ang{18;10;8}\) \\
    M64/NGC 4826 & 中間螺旋星系 & 黑眼星系 & \(\ang{;10;30}\times\ang{;5;20}\) & \(12^\text{h }56^\text{m }44^\text{s}\) & \(+\ang{21;40;58}\) \\
    M85/NGC 4382 & 透鏡狀星系 & - & \(\ang{;7;}\times\ang{;5;20}\) & \(12^\text{h }25^\text{m }23^\text{s}\) & \(+\ang{18;11;29}\) \\
    M88/NGC 4501 & 螺旋星系 & - & \(\ang{;8;40}\times\ang{;4;20}\) & \(12^\text{h }31^\text{m }59^\text{s}\) & \(+\ang{14;25;13}\) \\
    M91/NGC 4548 & 螺旋星系 & - & \(\ang{;5;24}\times\ang{;4;18}\) & \(12^\text{h }35^\text{m }26^\text{s}\) & \(+\ang{14;29;46}\) \\
    M98/NGC 4192 & 中間螺旋星系 & - & \(\ang{;11;}\times\ang{;2;40}\) & \(12^\text{h }13^\text{m }48^\text{s}\) & \(+\ang{14;54;01}\) \\
    M99/NGC 4254 & 螺旋星系 & \makecell{Coma Pinwheel\\Galaxy} & \(\ang{;5;}\) & \(12^\text{h }18^\text{m }50^\text{s}\) & \(+\ang{14;24;59}\) \\
    M100/NGC 4321 & 中間螺旋星系 & Mirror Galaxy & \(\ang{;6;}\) & \(12^\text{h }22^\text{m }54^\text{s}\) & \(+\ang{15;49;18}\) \\
\end{longtable}
\end{center}

\subsection{烏鴉座 Corvus}
\begin{enumerate}
    \item 最亮星：軫宿一，視星等2.59的衰老藍巨星
    \item 最亮的四顆恆星分別是軫宿一、軫宿三、軫宿二、軫宿四，在夜空構成顯著的四邊形
\end{enumerate}
烏鴉座中的深空天體：

\begin{center}
\begin{longtable}{c | c | c | c | c | c}
    % \hline
    編號 & 種類 & 別名 & 視尺寸 & RA & Dec \\
    \hline
    \endfirsthead

    % \hline
    % 編號 & 種類 & 別名 & 視尺寸 & 拍攝技巧 & 曝光時間 \\
    % \hline
    % \endhead

    % \hline
    % \multicolumn{6}{r}{接下頁} \\
    \endfoot

    % \hline
    \endlastfoot

    NGC 4038/NGC 4039 & 交互作用星系 & - & \(\ang{;6;}\times\ang{;4;}\) & \(12^\text{h }01^\text{m }53^\text{s}\) & \(-\ang{18;52;40}\) \\
\end{longtable}
\end{center}

\subsection{巨爵座 Crater}
\begin{enumerate}
    \item 最亮星：翼宿七，視星等3.56巨爵座沒有恆星亮度高於三等星
\end{enumerate}

\subsection{六分儀座 Sextans}
\begin{enumerate}
    \item 最亮的星（六分儀座\(\alpha\)）視星等4.49
\end{enumerate}

\subsection{長蛇座 Hydra}
\begin{enumerate}
    \item 最大的星座雖然面積很大（1303平方度），卻只包括一顆很亮的星星宿一，視星等1.98
    \item 傳說：在希臘神話中為太陽神阿波羅取水的烏鴉（烏鴉座）路上怠慢，而且取回了一條水蛇。阿波羅看穿了這個騙局，把烏鴉、水杯（巨爵座）和水蛇（長蛇座）扔到了天上。這些天區地帶在古代又被稱為「海」
    \item 又一說這是神話英雄赫拉克勒斯殺死的九頭蛇（海德拉），為前者的十二項英雄偉績之一
\end{enumerate}
長蛇座中的深空天體：

\begin{center}
\begin{longtable}{c | c | c | c | c | c}
    % \hline
    編號 & 種類 & 別名 & 視尺寸 & RA & Dec \\
    \hline
    \endfirsthead

    % \hline
    % 編號 & 種類 & 別名 & 視尺寸 & 拍攝技巧 & 曝光時間 \\
    % \hline
    % \endhead

    % \hline
    % \multicolumn{6}{r}{接下頁} \\
    \endfoot

    % \hline
    \endlastfoot

    M48/NGC 2548 & 疏散星團 & - & \(\ang{;28;12}\) & \(8^\text{h }13^\text{m }42^\text{s}\) & \(-\ang{5;45;2}\) \\
    M68/NGC 4590 & 球狀星團 & - & \(\ang{;6;48}\) & \(12^\text{h }39^\text{m }28^\text{s}\) & \(-\ang{26;44;35}\) \\
    M83/NGC 2536 & 螺旋星系 & 南風車星系 & \(\ang{;13;30}\) & \(13^\text{h }37^\text{m }00^\text{s}\) & \(-\ang{29;51;56}\)
\end{longtable}
\end{center}

\subsection{半人馬座 Centaurus}
\begin{enumerate}
    \item 最亮星：南門二（半人馬座\(\alpha\)星）是一顆三合星，距離太陽最近的恆星。
    \item 半人馬座\(\alpha\)星A：它與太陽擁有相同的化學構成並且擁有相同的亮度（0星等）。
    \item 半人馬座\(\alpha\)星C：又名比鄰星，為一顆紅矮星，亮度+12.4等
    \item 拆分為半人馬座與南十字座
\end{enumerate}
半人馬座中的深空天體：

\begin{center}
\begin{longtable}{c | c | c | c | c | c}
    % \hline
    編號 & 種類 & 別名 & 視尺寸 & RA & Dec \\
    \hline
    \endfirsthead

    % \hline
    % 編號 & 種類 & 別名 & 視尺寸 & 拍攝技巧 & 曝光時間 \\
    % \hline
    % \endhead

    % \hline
    % \multicolumn{6}{r}{接下頁} \\
    \endfoot

    % \hline
    \endlastfoot

    NGC 5139 & 球狀星團 & 半人馬座\(\omega\)星團 & \(\ang{;27;}\) & \(13^\text{h }26^\text{m }45^\text{s}\) & \(-\ang{47;28;36}\) \\
\end{longtable}
\end{center}

\subsection{南十字座 Crux}
\begin{enumerate}
    \item 可用來標示天球南極
    \item 最小的星座
    \item 十字架一是一對雙星，主星是紅色的巨星，視星等1.6等，伴星視星等6.5等
    \item 此外相對於南十字，北天也有北十字（天鵝座的中心部分），在其西方另有假十字
    \item 十字架一（\(\gamma\)）、十字架二（\(\alpha\)）、十字架三（\(\beta\)）、十字架四（\(\delta\)）
\end{enumerate}
南十字座中的深空天體：

\begin{center}
\begin{longtable}{c | c | c | c | c | c}
    % \hline
    編號 & 種類 & 別名 & 視尺寸 & RA & Dec \\
    \hline
    \endfirsthead

    % \hline
    % 編號 & 種類 & 別名 & 視尺寸 & 拍攝技巧 & 曝光時間 \\
    % \hline
    % \endhead

    % \hline
    % \multicolumn{6}{r}{接下頁} \\
    \endfoot

    % \hline
    \endlastfoot

    C99 & 暗星雲 & 媒袋星雲 & \(\ang{7;;}\times\ang{5;;}\) & \(12^\text{h }50^\text{m }00^\text{s}\) & \(-\ang{62;30;00}\) \\
    NGC 4755 & 疏散星團 & 珠寶盒星團 & \(\ang{;7;48}\) & \(12^\text{h }53^\text{m }36^\text{s}\) & \(-\ang{60;21;23}\)
\end{longtable}
\end{center}

\subsection{羅盤座 Pyxis}
\begin{enumerate}
    \item 星座中最亮的三顆星分別是天狗五、天狗四和天狗六且基本排成直線。3.68的天狗五是羅盤座最亮恆星
\end{enumerate}

\subsection{唧筒座 Antlia}
\begin{enumerate}
    \item 最亮的恆星是疑似變星的橘色巨星唧筒座\(\alpha\)，視星等為4.22至4.29
    \item 拉丁語義「泵」；代表氣泵
\end{enumerate}

\subsection{圓規座 Circinus}
\begin{enumerate}
    \item 只有一顆亮於視星等4的恆星。圓規座α是一顆視星等為3.19的白色
    \item 在半人馬座下，無法觀測
\end{enumerate}
圓規座中的深空天體：

\begin{center}
\begin{longtable}{c | c | c | c | c | c}
    % \hline
    編號 & 種類 & 別名 & 視尺寸 & RA & Dec \\
    \hline
    \endfirsthead

    % \hline
    % 編號 & 種類 & 別名 & 視尺寸 & 拍攝技巧 & 曝光時間 \\
    % \hline
    % \endhead

    % \hline
    % \multicolumn{6}{r}{接下頁} \\
    \endfoot

    % \hline
    \endlastfoot

    ESO 97-G13 & 不規則狀星系 & 圓規座星系 & \(\ang{;7;}\times\ang{;3;}\) & \(14^\text{h }13^\text{m }10^\text{s}\) & \(-\ang{65;20;21}\) \\
\end{longtable}
\end{center}